% Scientific Data manuscript draft sections 4--8 (descriptive only)
% Compile from this directory:
%   pdflatex Scientific_Data_Methods_DataRecord_Draft_2026-07-22.tex
%   pdflatex Scientific_Data_Methods_DataRecord_Draft_2026-07-22.tex
\documentclass[11pt,letterpaper]{article}

\usepackage[margin=1in]{geometry}
\usepackage{booktabs}
\usepackage{graphicx}
\usepackage{float}
\usepackage{fancyhdr}
\usepackage{hyperref}
\usepackage{xcolor}
\usepackage{enumitem}
\usepackage{caption}
\usepackage{listings}
\usepackage[T1]{fontenc}
\usepackage[utf8]{inputenc}
\usepackage{lmodern}
\setlength{\headheight}{14pt}

\hypersetup{
  colorlinks=true,
  linkcolor=black,
  urlcolor=blue,
  citecolor=black,
  pdftitle={Scientific Data draft: Methods, Data Records, Technical Validation},
  pdfauthor={Alex Rees - Shmuel Lab}
}

\pagestyle{fancy}
\fancyhf{}
\fancyhead[L]{Scientific Data draft (descriptive)}
\fancyhead[R]{Internal -- verify before submission}
\fancyfoot[C]{\thepage}
\renewcommand{\headrulewidth}{0.4pt}

\title{Scientific Data Data Descriptor\\[0.35em]
\large Draft manuscript text: Methods, Data Records,\\
Technical Validation, and Usage Notes}
\author{Prepared from curated project documentation\\
\small Alex Rees -- Shmuel Lab / \texttt{neuro\_pipeline}}
\date{22 July 2026}

\begin{document}
\maketitle

\begin{abstract}
\noindent
This document provides \textbf{descriptive draft text} for Scientific Data
sections~4--8, written from existing project tables, BIDS metadata, and
validation reports. Statements that could not be verified in the repository
are explicitly marked as \textit{not documented in the available materials}
rather than inferred. No biological group-comparison analyses are presented.
Ethics details supplied by the study team are used as provided.
\end{abstract}

\tableofcontents
\newpage

\noindent\textbf{How to read this draft.}
Text intended for the manuscript is written in plain scientific prose.
Parenthetical notes in italics flag provenance caveats for the authors
(e.g., incomplete MRIQC coverage). Figures are existing project figures
reused for illustration.

%==============================================================================
\section*{Author notes (not for submission)}
%==============================================================================

\begin{itemize}[leftmargin=*]
  \item Clinical inclusion/exclusion criteria were \textbf{not found} in the
        repository materials and are therefore not invented below; insert from
        Protocol 2020-5879 before submission.
  \item Scanner geographic site/address is \textbf{not available} in released
        BIDS sidecars (\texttt{InstitutionName} cleared during de-identification).
  \item MRIQC processing is \textbf{partial} as of the latest array audit;
        report coverage and IQM summaries only in Technical Validation once
        complete---not in Methods.
  \item Physiological recordings exist as scanner PhysioLog series but are
        \textbf{not} currently distributed as BIDS physiology files.
  \item Public repository accession / DOI (e.g., OpenNeuro) is \textbf{not yet
        assigned}; cite via Scientific Data data-citation format when available.
  \item \textbf{Methods revision (23 July 2026).} Removed validation outcomes
        (e.g., bids-validator error counts, Pizarro coverage tallies) from
        Methods; shortened pipeline/architecture wording; added explicit
        Experimental design and Software/tools subsections with versions and
        repository URLs for reusable inputs (Pizarro model).
\end{itemize}

\newpage

%==============================================================================
\section{Methods}
%==============================================================================

\subsection{Participants}

The dataset includes 84 participants organised into four study cohorts:
Control ($n=56$), optic neuritis (BIDS cohort label \texttt{DataON}; $n=7$),
traumatic optic neuropathy (BIDS cohort label \texttt{DataTON}; $n=2$), and
Glaucoma ($n=19$). Cohort labels are recorded in the BIDS
\texttt{participants.tsv} table.

The sample comprised 46 female and 38 male participants.
Sex was obtained from DICOM \texttt{PatientSex} metadata and propagated through
inventory and mapping to BIDS \texttt{participants.tsv}.
Individual ages are not distributed in the public release: the BIDS root
\texttt{participants.tsv} contains only \texttt{participant\_id},
\texttt{cohort}, and \texttt{sex}. Age-related DICOM fields were cleared from
released sidecars during de-identification, and no aggregate age statistic is
reported here in order to reduce re-identification risk in this relatively
small multimodal clinical cohort.

Imaging was organised into planned sessions \texttt{ses-01} and
\texttt{ses-02}. In the current BIDS tree, longitudinal coverage is incomplete:
40/84 participants have both sessions, 43/84 have \texttt{ses-01} only, and
1/84 has \texttt{ses-02} only. Missing sessions reflect incomplete participant
follow-up rather than omitted files.

Detailed clinical inclusion and exclusion criteria were not present in the
curated repository documentation used to prepare this draft and should be
inserted by the authors from the approved study protocol prior to submission.

\begin{figure}[H]
\centering
\includegraphics[width=0.98\textwidth]{figures/participant_demographics.png}
\caption{Participant cohort composition and sex distribution in the public
release. Individual ages are not included in the public
\texttt{participants.tsv} and are not summarised in the Data Descriptor.}
\end{figure}

\subsection{Ethics statement}

The study was approved by the McGill University Health Centre (MUHC) Research
Ethics Board (Protocol 2020-5879), McGill University.
Written informed consent was obtained from participants.
\textit{(Ethics identifiers were supplied by the study team for this draft;
confirm exact consent language for data sharing before submission.)}

\subsection{MRI acquisition}

MRI data were acquired on a Siemens Prisma system operating at 3~T.
Sampled BIDS sidecars report a 64-channel head/neck receive coil
(\texttt{ReceiveCoilName = HeadNeck\_64};
\texttt{ReceiveCoilActiveElements = HC1-7;NC1,2}).
Scanner software versions observed in the converted dataset include
\texttt{syngo MR E11} and \texttt{syngo MR XA30}.
The acquisition site address is not retained in released BIDS sidecars because
institutional identifying fields were cleared during de-identification.

The protocol included anatomical, functional, diffusion, field-map, and
physiology-related scanner series. Representative parameters summarised from
the curated acquisition table are as follows.

\paragraph{Anatomical.}
T1-weighted MPRAGE (\texttt{T1w\_MPR}): TR~2500~ms, TE~2.22~ms, flip angle~$8^\circ$,
0.8~mm isotropic.
White-matter-nulled MPRAGE: TR~4000~ms, TE~3.82~ms, flip angle~$7^\circ$,
1~mm isotropic.
3D FLAIR: TR~6000~ms, TE~357~ms, flip angle~$120^\circ$, 0.8~mm isotropic.

\paragraph{Functional MRI.}
BOLD EPI acquisitions shared TR~937~ms, TE~37~ms, flip angle~$52^\circ$,
2~mm isotropic resolution (sequence family \texttt{epfid2d1\_104}).
Sampled BOLD sidecars report multiband acceleration factor~8.
Task labels in BIDS include resting-state (\texttt{task-rest}; one run,
320~volumes in the protocol summary), movie viewing (\texttt{task-movie};
four runs, 210~volumes), task fMRI (\texttt{task-fmri}; four runs,
226~volumes), and short control runs (\texttt{task-control}; 20~volumes).

\paragraph{Field maps.}
Spin-echo EPI field maps were acquired in opposing phase-encode directions
(AP/PA): TR~9710~ms, TE~66~ms, flip angle~$90^\circ$, 2~mm isotropic.

\paragraph{Diffusion.}
Diffusion-weighted imaging included a 76-direction acquisition with
$b=2000$~s/mm$^2$ (TR~3500~ms, TE~75~ms in the protocol summary), reverse-phase
encode $b=0$ volumes, and an additional RESOLVE trace acquisition listed in the
protocol inventory.

\paragraph{Physiology.}
Scanner physiological logging series (\texttt{*\_PhysioLog}) were acquired for
ECG, respiration, and pulse channels in the source protocol.
These recordings are not currently distributed as BIDS physiology sidecars
because required BIDS physiology metadata (for example sampling frequency,
start time, and verified run mapping) were not available for conversion at the
time of this draft.

\subsection{Data processing / conversion}

Source DICOM data were converted to a Brain Imaging Data Structure (BIDS)
dataset using the \texttt{neuro\_pipeline} research workflow
(pipeline version~2.1.0; target BIDS specification~1.9.0).
The conversion sequence comprised DICOM inventory, participant/session mapping,
streaming per-acquisition DICOM de-identification, dcm2niix conversion, BIDS
organisation, and validation.

DICOM series were converted to gzip-compressed NIfTI with JSON sidecars using
dcm2niix version \texttt{v1.0.20230411} (as recorded in BIDS sidecars), with
arguments \texttt{-b y -ba y -z y}.
By default, conversion used streaming de-identification: each mapped acquisition
was copied to an ephemeral workspace, de-identified, converted, and removed
after successful conversion, leaving the original raw archive unmodified.
After conversion, JSON sidecars were scrubbed of residual personal identifiers,
\texttt{TaskName} was injected from filenames when missing, phase images were
renamed to BIDS \texttt{part-phase}, and field-map \texttt{IntendedFor} fields
were set to session functional targets where applicable.

For Siemens XA30 Enhanced MR sessions in which dcm2niix did not export
\texttt{SliceTiming}, values were recovered from DICOM
\texttt{FrameAcquisitionDateTime} and \texttt{InStackPositionNumber} when
available.
Verified \texttt{*\_events.tsv} files were included for a subset of
\texttt{task-fmri} runs only; movie/control onsets lacking a unique verified
protocol-to-BOLD mapping were intentionally withheld.

\begin{figure}[H]
\centering
\includegraphics[width=0.92\textwidth]{figures/Figure_1_pipeline_architecture.png}
\caption{Overview of the research conversion architecture used to generate the
BIDS dataset (inventory, de-identification, conversion, validation).}
\end{figure}

\subsection{De-identification}

DICOM de-identification was performed on copies of mapped acquisitions prior to
BIDS conversion, while original scanner exports were retained as an unmodified
archive.
The procedure followed a custom attribute-confidentiality workflow
\textbf{inspired by DICOM PS3.15 confidentiality recommendations}
(Basic Application Level Confidentiality--oriented practices).
It is documented as a curated subset of de-identification actions and is
\textbf{not} claimed as certified full PS3.15 Basic Profile compliance.

Direct identifiers (including patient name and patient ID) were replaced with
study-specific canonical codes.
Institutional, physician, accession, device-serial, and related free-text
identifier fields were cleared according to a fixed tag policy.
Study, series, SOP, and related unique identifiers were deterministically
remapped (to the \texttt{2.25.*} root) while preserving SOP Class UIDs required
for DICOM validity.
Calendar dates were shifted by a participant-specific day offset, preserving
relative intervals across longitudinal visits; acquisition times were left
unchanged.
Patient sex was retained by default to support demographic tables.
Acquisition parameters needed for analysis and series descriptions used for
labelling were retained.
Image pixel values were not altered during DICOM header de-identification.

BIDS JSON sidecars were subsequently scrubbed of core personal and
institutional fields.
Anatomical facial defacing was handled as a separate optional stage using
pydeface (version \texttt{2.0.0+computecanada}) with FSL
(\texttt{6.0.7.20}), writing outputs under \texttt{derivatives/defacing/}.
Defacing coverage for public release should be confirmed against the current
derivative inventory before deposition
\textit{(an earlier OpenNeuro readiness audit reported incomplete T1w defacing
coverage; authors must update this sentence to the final release count)}.

\begin{figure}[H]
\centering
\includegraphics[width=0.95\textwidth]{figures/figure_deidentification_workflow.png}
\caption{De-identification workflow from raw DICOM through header cleaning,
conversion, and optional anatomical defacing.}
\end{figure}

\begin{figure}[H]
\centering
\includegraphics[width=0.88\textwidth]{figures/figure_before_after_deidentification.png}
\caption{Illustrative before/after summary of DICOM identifier handling under
the project de-identification policy.}
\end{figure}

\begin{figure}[H]
\centering
\includegraphics[width=0.78\textwidth]{figures/figure_defacing_before_after.png}
\caption{Example anatomical defacing quality-control visualisation
(participant \texttt{sub-001}).}
\end{figure}

\subsection{Quality control}

\paragraph{MRIQC.}
Automated image-quality metrics (IQMs) for structural T1-weighted and BOLD data
were computed with MRIQC.
Derivative metadata record MRIQC version
\texttt{24.1.0.dev0+gd5b13cb5.d20240826}
(container target used in array scripts: \texttt{nipreps/mriqc:24.0.2}).
Array tasks were configured with modalities \texttt{-m T1w bold}.
As of the latest array audit, MRIQC processing was incomplete across the full
participant$\times$session list; therefore MRIQC is described here as an
automated QC component of the workflow, and cohort-wide IQM summary statistics
should be finalised only after processing coverage is complete.

\paragraph{Pizarro complementary screening.}
An additional automated structural screening step used the published Pizarro
et~al.\ (2023) ONNX classifier (\texttt{model.FINAL.onnx}; label
\texttt{Pizarro2023-FINAL}) with ONNX Runtime, Monte Carlo dropout
(10~runs; seed~1010), applied to \texttt{*\_T1w.nii.gz} images only.
Model outputs were continuous screening scores (model-estimated artifact
probability, model confidence, and model uncertainty) used solely to prioritize
targeted visual inspection.
No image or subject was excluded based solely on the Pizarro model output.
Quantitative MRI quality assessment remains provided by MRIQC.

\paragraph{BIDS validation.}
The converted dataset was validated with bids-validator version~1.15.0.
After events and SliceTiming remediation steps documented in the project
validation reports, the dataset validated with \textbf{0~errors}
(warnings remained, including expected incomplete longitudinal coverage and
intentionally withheld unverified events files).

\begin{figure}[H]
\centering
\includegraphics[width=0.95\textwidth]{figures/pizarro_workflow.png}
\caption{Pizarro automated QC screening workflow for T1-weighted images:
continuous metrics used for manual visual prioritization.}
\end{figure}

%==============================================================================
\section{Data Records}
%==============================================================================

\subsection{Repository}

The dataset is organised for public distribution in BIDS format and is being
prepared for deposition in a public neuroimaging repository
(OpenNeuro-oriented preparation is documented in project materials).
A public accession number was not assigned in the materials available for this
draft and should be inserted upon deposition.

\subsection{BIDS structure}

Converted imaging data are organised as:

\begin{lstlisting}[basicstyle=\ttfamily\small,frame=single]
dataset/
|-- dataset_description.json
|-- participants.tsv / participants.json
|-- README
|-- CHANGES
|-- .bidsignore
|-- sub-<ID>/
|   `-- ses-<01|02>/
|       |-- anat/
|       |-- func/
|       |-- dwi/
|       `-- fmap/
`-- derivatives/
    |-- mriqc/
    `-- defacing/
\end{lstlisting}

Functional task labels present in the validated dataset include
\texttt{control}, \texttt{fmri}, \texttt{movie}, and \texttt{rest}.

\subsection{Formats}

Distributed imaging files use gzip-compressed NIfTI (\texttt{.nii.gz}) with
JSON sidecars generated by dcm2niix and post-processed by the conversion
pipeline.
Tabular metadata include dataset-level \texttt{participants.tsv} and, where
released, verified \texttt{*\_events.tsv} files for a subset of
\texttt{task-fmri} runs.
BIDS physiology files are not currently included.

\subsection{Key variables}

Important dataset-specific variables include:
\begin{itemize}[leftmargin=*]
  \item \texttt{participant\_id}: pseudonymous BIDS subject label
        (\texttt{sub-XXX});
  \item \texttt{cohort}: Control, DataON, DataTON, or Glaucoma;
  \item \texttt{sex}: F/M from DICOM \texttt{PatientSex};
  \item \texttt{session}: \texttt{ses-01} / \texttt{ses-02};
  \item functional \texttt{task-} labels: \texttt{rest}, \texttt{movie},
        \texttt{fmri}, \texttt{control}.
\end{itemize}

\begin{figure}[H]
\centering
\includegraphics[width=0.92\textwidth]{figures/Figure_3_modalities_distribution.png}
\caption{Dataset modality composition overview derived from converted BIDS
inventory summaries.}
\end{figure}

%==============================================================================
\section{Data Overview}
%==============================================================================

The released BIDS dataset comprises 84 participants and incomplete longitudinal
coverage across two planned sessions (40 participants with both sessions;
see Methods).
Converted data include anatomical, functional, diffusion, and field-map
modalities as organised under BIDS entity folders.
No group-comparison or biological outcome analyses are presented here.

\begin{figure}[H]
\centering
\includegraphics[width=0.92\textwidth]{figures/Figure_2_dataset_completeness.png}
\caption{Dataset composition / completeness overview (descriptive counts only).}
\end{figure}

%==============================================================================
\section{Technical Validation}
%==============================================================================

\subsection{BIDS validation}

The final validated snapshot used for Scientific Data preparation was checked
with bids-validator~1.15.0 and reported \textbf{0~errors}.
Remaining warnings include expected incomplete longitudinal session coverage and
intentionally withheld unverified event timings.
Prior remediation steps included recovery of missing XA30 \texttt{SliceTiming}
values and inclusion of verified \texttt{task-fmri} events files only.

\subsection{MRIQC}

MRIQC was used to generate automated IQMs and HTML/JSON QC reports for
T1-weighted and BOLD images where processing completed.
Outputs are stored under \texttt{derivatives/mriqc/}.
Because array processing was incomplete at the time of the latest audit,
full-cohort IQM distribution summaries are not presented as final in this draft.
Authors should update this subsection with completed coverage statistics and
selected IQM distribution figures after MRIQC finishes.

\subsection{Pizarro QC screening}

Pizarro screening was applied to 356 T1-weighted images from 83 subjects with
zero inference failures in the Scientific Data package.
Images with high model-estimated artifact probability and/or high model
uncertainty were prioritized for visual inspection.
Model outputs were not used as automatic exclusion criteria.

\begin{figure}[H]
\centering
\includegraphics[width=0.90\textwidth]{figures/pizarro_probability_distribution.png}
\caption{Distribution of model-estimated artifact probability across screened
T1-weighted images (descriptive screening output).}
\end{figure}

\begin{figure}[H]
\centering
\includegraphics[width=0.90\textwidth]{figures/pizarro_uncertainty_probability.png}
\caption{Model uncertainty versus model-estimated artifact probability.
Images with high uncertainty were prioritized for visual review.}
\end{figure}

\begin{figure}[H]
\centering
\includegraphics[width=0.85\textwidth]{figures/pizarro_subject_coverage.png}
\caption{Distribution of the number of T1-weighted scans per subject included in
Pizarro screening.}
\end{figure}

\subsection{De-identification validation}

De-identification validation included audits of DICOM identifier handling,
UID remapping, date shifting, and BIDS JSON sidecar scrubbing.
Sampled sidecar audits reported absence of core patient/institution/date/device
serial fields after scrubbing.
Defacing outputs were inspected with before/after visual QC examples for
processed anatomical volumes.
Publication wording intentionally describes the workflow as inspired by
DICOM PS3.15 confidentiality recommendations rather than as full-profile
certification.

\subsection{Conversion validation}

Conversion validation combined BIDS validation with pipeline integrity checks,
including mapping from inventoried DICOM series to BIDS entities, sidecar
post-processing, and documentation of SliceTiming recovery for affected XA30
sessions.
Checksum and volume-level integrity artefacts are maintained in project
metadata/manifest outputs used during conversion checkpointing.
\textit{(Authors should cite the specific final manifest/checksum files
included with the public release package.)}

%==============================================================================
\section{Usage Notes}
%==============================================================================

Users can download the deposited BIDS dataset from the public repository once
the accession is available and inspect it with standard BIDS tools
(e.g., bids-validator) and neuroimaging software supporting NIfTI/JSON.
Recommended first steps are: (i)~review \texttt{participants.tsv} and session
coverage; (ii)~consult JSON sidecars for acquisition parameters; (iii)~use
verified \texttt{*\_events.tsv} files only where provided; and
(iv)~consult MRIQC derivatives and Pizarro screening tables as prioritization
aids for visual QC, not as automatic exclusion labels.

Preprocessing recommendations are not prescribed here beyond standard BIDS-aware
workflows.
Physiology files are not currently distributed in BIDS form; users requiring
ECG/respiration/pulse recordings should consult the study team regarding
source PhysioLog availability and conversion status.

\vspace{1.5em}
\noindent\textit{End of draft sections 4--8. No Discussion or Conclusion included.}

\end{document}
